\subsection{Paging Stage}

Paging is enabled by (:ref:base.control_registers.PTCR.PE:) and translates the (:term:base.terms.linear_address:) produced by (:term:base.terms.segment_pre_translation:). Before
the walk can produce a translation, the (:term:base.terms.linear_address:) must be canonical for the translation-table format selected by
(:ref:base.control_registers.PTCR.TT:). A noncanonical address raises (:ref:base.events.EXCEPTION.NONCANONICAL_ADDRESS:). With paging disabled, the memory system uses
the (:term:base.terms.linear_address:) directly and does not interpret (:ref:base.control_registers.PTCR.TT:) or (:ref:base.control_registers.PTCR.ROOT_PAGE:).

Every table page is one 16-KiB-aligned physical frame, and every page-table entry is one naturally aligned 64-bit descriptor. L4, L3, and L2 use
an 11-bit address (:term:base.terms.field:) to select one of 2048 entries. L1 uses a nine-bit address (:term:base.terms.field:) to select one of 512 entries in bytes
\texttt{0x0000..0x0fff}; the walker does not access bytes \texttt{0x1000..0x3fff} of an L1 table page. For a table-page base address $B$ and
index $i$, the selected entry is always at $B+8i$. PTEs are a dense array of eight-byte descriptors. The architecture does not define packed
L1 tables, sub-page table pointers, or table fragments. The trailing L1 bytes are not reserved storage attached to individual PTEs.

The walker starts at the physical table page named by (:ref:base.control_registers.PTCR.ROOT_PAGE:) for the selected format. At each active level,
a present selected entry with (:ref:base.memory.translation.PTE.T:)\texttt{=1} points to the next 16-KiB table page; this form is valid only above L1.
A selected entry with (:ref:base.memory.translation.PTE.T:)\texttt{=0} terminates the walk as a byte-addressed leaf. Define the number of low
linear-address bits supplied by a leaf at level $L$ as $o(L1)=14$, $o(L2)=23$, $o(L3)=34$, and $o(L4)=45$. The leaf
(:ref:base.memory.translation.PTE.PFN:) supplies a naturally aligned physical base, and the final byte address is:

\[
  \mathit{byte\_address} = (\mathit{leaf\_PFN} \ll 14) \mathbin{|}
  \mathit{linear}[o(L)-1:0].
\]

The resulting L1 through L4 leaf sizes are 16 KiB, 8 MiB, 16 GiB, and 32 TiB. Bit 7 is (:ref:base.memory.translation.PTE.A:) in both descriptor
layouts and does not participate in address formation.

The walker intersects table-pointer (:ref:base.memory.translation.PTE.R:)/(:ref:base.memory.translation.PTE.W:)/(:ref:base.memory.translation.PTE.X:)/(:ref:base.memory.translation.PTE.U:) (:term:base.terms.field|plural:) while descending. The
terminating leaf supplies decoded (:ref:base.memory.translation.PTE.AM:) permissions, (:ref:base.memory.translation.PTE.U:), the final (:ref:base.memory.translation.PTE.PFN:),
(:ref:base.memory.translation.PTE.G:), (:ref:base.memory.translation.PTE.CP:), and the Normal/MMIO class. The detailed entry-validity,
permission, accessed, and dirty rules are specified under Page-Table Entry Format. At each consumed level, result
priority is not-present first, structural validity second, and effective permission third; the walker resolves
those results before consuming a lower level.

\begin{BedrockListedFormatDiagram}{Linear-Address Paging Fields}
\BedrockFormatRowRange{LA45 linear[63:0]}{63}{0}{%
\BedrockFormatField{sign extension}{19}
\BedrockFormatField{L3}{11}
\BedrockFormatField{L2}{11}
\BedrockFormatField{L1}{9}
\BedrockFormatField{offset}{14}
}
\BedrockFormatRowRange{LA56 linear[63:0]}{63}{0}{%
\BedrockFormatField{SEXT}{8}
\BedrockFormatField{L4}{11}
\BedrockFormatField{L3}{11}
\BedrockFormatField{L2}{11}
\BedrockFormatField{L1}{9}
\BedrockFormatField{offset}{14}
}
\end{BedrockListedFormatDiagram}

\clearpage
